% Options for packages loaded elsewhere
\PassOptionsToPackage{unicode}{hyperref}
\PassOptionsToPackage{hyphens}{url}
%
\documentclass[
]{article}
\usepackage{amsmath,amssymb}
\usepackage{iftex}
\ifPDFTeX
  \usepackage[T1]{fontenc}
  \usepackage[utf8]{inputenc}
  \usepackage{textcomp} % provide euro and other symbols
\else % if luatex or xetex
  \usepackage{unicode-math} % this also loads fontspec
  \defaultfontfeatures{Scale=MatchLowercase}
  \defaultfontfeatures[\rmfamily]{Ligatures=TeX,Scale=1}
\fi
\usepackage{lmodern}
\ifPDFTeX\else
  % xetex/luatex font selection
\fi
% Use upquote if available, for straight quotes in verbatim environments
\IfFileExists{upquote.sty}{\usepackage{upquote}}{}
\IfFileExists{microtype.sty}{% use microtype if available
  \usepackage[]{microtype}
  \UseMicrotypeSet[protrusion]{basicmath} % disable protrusion for tt fonts
}{}
\makeatletter
\@ifundefined{KOMAClassName}{% if non-KOMA class
  \IfFileExists{parskip.sty}{%
    \usepackage{parskip}
  }{% else
    \setlength{\parindent}{0pt}
    \setlength{\parskip}{6pt plus 2pt minus 1pt}}
}{% if KOMA class
  \KOMAoptions{parskip=half}}
\makeatother
\usepackage{xcolor}
\usepackage[margin=2.5cm]{geometry}
\usepackage{longtable,booktabs,array}
\usepackage{calc} % for calculating minipage widths
% Correct order of tables after \paragraph or \subparagraph
\usepackage{etoolbox}
\makeatletter
\patchcmd\longtable{\par}{\if@noskipsec\mbox{}\fi\par}{}{}
\makeatother
% Allow footnotes in longtable head/foot
\IfFileExists{footnotehyper.sty}{\usepackage{footnotehyper}}{\usepackage{footnote}}
\makesavenoteenv{longtable}
\setlength{\emergencystretch}{3em} % prevent overfull lines
\providecommand{\tightlist}{%
  \setlength{\itemsep}{0pt}\setlength{\parskip}{0pt}}
\setcounter{secnumdepth}{-\maxdimen} % remove section numbering
\ifLuaTeX
  \usepackage{selnolig}  % disable illegal ligatures
\fi
\IfFileExists{bookmark.sty}{\usepackage{bookmark}}{\usepackage{hyperref}}
\IfFileExists{xurl.sty}{\usepackage{xurl}}{} % add URL line breaks if available
\urlstyle{same}
\hypersetup{
  hidelinks,
  pdfcreator={LaTeX via pandoc}}

\author{}
\date{}

\begin{document}

\hypertarget{linstabilituxe0-essenziale-del-vuoto-di-yangmills-riproduzione-machine-checked-del-programma-di-preparata}{%
\section{L'instabilità essenziale del vuoto di Yang--Mills: riproduzione
machine-checked del programma di
Preparata}\label{linstabilituxe0-essenziale-del-vuoto-di-yangmills-riproduzione-machine-checked-del-programma-di-preparata}}

Versione 1.0 --- 11 agosto 2026

Giuliano (Vantasner AG), per Luca Gamberale

Rapporto tecnico con verifica automatica (substrate-framework, tracker
\#44)

\hypertarget{sommario}{%
\subsection{Sommario}\label{sommario}}

Questo rapporto riproduce, con oracoli macchina (SymPy/SciPy/mpmath, 50
test automatici), il programma di Giuliano Preparata sull'instabilità
del vuoto perturbativo di Yang--Mills: il potenziale efficace a un loop
nel fondo cromomagnetico costante (vuoto di Savvidy), il potenziale a
due loop in funzione di \(b/\Lambda^2\) (con \(b = gH\)), l'analisi di
stabilità rispetto alle fluttuazioni, la decomposizione per colori e per
modi trasversi, e una configurazione classica migliore di quella di
Savvidy (il vuoto ``spaghetti'' di Ambjørn--Olesen) con la
riquantizzazione attorno ad essa. Ogni affermazione quantitativa è
prodotta da codice eseguibile nel repository
\texttt{vantasnerdan/substrate-framework} (issue \#44--\#50, PR
\#51--\#54) o citata esplicitamente dalla letteratura. Risultato
principale: il risultato a un loop di Savvidy si riproduce esattamente;
il risultato a due loop pubblicato (Bordag--Skalozub 2022) contiene tre
errori tipografici/sostanziali che documentiamo con valori esatti; la
configurazione di Savvidy è instabile alle fluttuazioni a ogni ordine
perturbativo accessibile, e il condensato di tubi di flusso (reticolo
triangolare, parametro di Abrikosov \(\beta = 1.159595\)) è
energeticamente preferito.

\hypertarget{fonti-e-accessibilituxe0}{%
\subsection{1. Fonti e accessibilità}\label{fonti-e-accessibilituxe0}}

Dichiariamo esattamente cosa è stato verificato e da dove.

\begin{itemize}
\tightlist
\item
  G. Preparata, \emph{Essential Quantum Instability of the Perturbative
  Yang--Mills Vacuum}, Nuovo Cim. A \textbf{96} (1986) 366, DOI
  10.1007/BF02833896 --- \textbf{testo completo non accessibile}
  (paywall Springer; l'archivio Sapienza è solo fisico). Del paper
  centrale citiamo solo il livello dell'abstract: un argomento
  variazionale con cut-off UV per cui
  \(\Delta E \sim -b\,\Lambda_{UV}^4\) con \(b\) finito per \(g \to 0\)
  --- instabilità ``essenziale'', cioè non riparabile perturbativamente.
  Se ci fornisci il PDF, la verifica equazione per equazione entra nel
  workflow.
\item
  G.K. Savvidy, Phys. Lett. B \textbf{71} (1977) 133 --- il risultato a
  un loop, ripreso dalle sue riesposizioni aperte (EPJC \textbf{80}
  (2020) 165, arXiv:1910.00654; la rassegna EPJC 2026). Verificato
  interamente.
\item
  N.K. Nielsen, P. Olesen, Nucl. Phys. B \textbf{144} (1978) 376 ---
  paywalled; il modo instabile e la parte immaginaria sono stati
  verificati indipendentemente dal nostro calcolo (zeta dell'operatore e
  integrale di modo).
\item
  M. Bordag, V. Skalozub, EPJC \textbf{82} (2022) 390, arXiv:2112.01043
  --- \textbf{aperto}; è la fonte del due loop a \(T=0\) (eq. 57--58).
  Verificato; tre discrepanze documentate (Sezione 5).
\item
  M. Bordag, \emph{Symmetry} \textbf{15} (2023) 1137 --- rassegna
  aperta; fonte della costruzione del reticolo di tubi di flusso (eq.
  146--152).
\item
  P. Cea, arXiv:2311.14791 --- aperto; settori tachiometrici \(SU(3)\).
\end{itemize}

Metodo: nessun passaggio ``a mano''. Ogni coefficiente è prodotto da
codice; ogni check ha una mutazione che lo rompe. I moduli:
\texttt{chromomagnetic\_background.py} (21 test),
\texttt{chromomagnetic\_two\_loop.py} (13 test),
\texttt{chromomagnetic\_sectors.py} (10 test),
\texttt{spaghetti\_vacuum.py} (8 test).

\hypertarget{il-potenziale-a-un-loop-savvidy-verificato}{%
\subsection{2. Il potenziale a un loop (Savvidy),
verificato}\label{il-potenziale-a-un-loop-savvidy-verificato}}

Convenzioni congelate: segnatura mostly-plus, gauge di fondo di Feynman
(\(\xi = 1\)), \(SU(2)\) con fondo in una direzione di Cartan
(\(a = 3\)), variabile covariante \(b = gH\).

\textbf{Spettro delle fluttuazioni} (derivato dall'operatore
discretizzato, non asserito): vettore carico
\[E^2 = p_z^2 + b(2n+1) - 2b\,s_3, \qquad s_3 \in \{+1, 0, 0, -1\},\]
fantasma carico \(E^2 = p_z^2 + b(2n+1)\). Il livello \(n=0\),
\(s_3 = +1\) dà \(E^2 = p_z^2 - b\): il \textbf{tachione di
Nielsen--Olesen}.

\textbf{Cancellazione fantasma}: la traccia di spin \(2\cosh(2bs) + 2\)
meno il fantasma (\(-2\)) lascia \(2\cosh(2bs)\); il coefficiente di
Seeley--DeWitt è esatto:
\[\frac{b}{\sinh(bs)}\,2\cosh(2bs) = \frac{2}{s} + \frac{11}{3}b^2 s + \frac{127}{180}b^4 s^3 + \cdots\]

\textbf{Potenziale a un loop} (MS-bar):
\[V_1(b) = \frac{11}{48\pi^2}\,b^2\left(\ln\frac{b}{\mu^2} - \frac{1}{2}\right).\]
Il coefficiente \(11/(48\pi^2)\) è fissato a \(10^{-10}\) dalla rotta
zeta dell'operatore e a \(10^{-2}\) dalla rotta proper-time con cut-off,
indipendenti. Il minimo:
\[\ln\frac{b_{\min}}{\mu^2} = -\frac{24\pi^2}{11 g^2}, \qquad V_{\min} = -\frac{11\mu^4}{96\pi^2}\,e^{-48\pi^2/11g^2} < 0,\]
identico alla formula pubblicata da Savvidy. Con il running a un loop
(\(b_0 = 22/3\)): \(b_{\min} = \Lambda^2\) \textbf{esattamente} ---
contenuto indipendente dallo schema.

\hypertarget{stabilituxe0-rispetto-alle-fluttuazioni}{%
\subsection{3. Stabilità rispetto alle
fluttuazioni}\label{stabilituxe0-rispetto-alle-fluttuazioni}}

Il verdetto è \textbf{negativo} e lo enunciamo con precisione:

\begin{itemize}
\tightlist
\item
  Per ogni \(b > 0\) esiste il modo tachiometrico
  \(E^2 = p_z^2 - b < 0\) per \(p_z^2 < b\).
\item
  Parte immaginaria a un loop, derivata dall'integrale di modo e
  indipendentemente dalla zeta dell'operatore:
  \[|\,\mathrm{Im}\,V_1| = \frac{b^2}{8\pi}.\]
\item
  A due loop la parte immaginaria persiste: dal quadrato complesso
  \((3g^2/2)B_2^2\) con \(B_2 = -b(\ln 2 - i\pi)/(16\pi^2)\) si ottiene
  \(\mathrm{Im}\,V_2 = -3g^2 b^2 \ln 2/(256\pi^3) \neq 0\).
\item
  La risommazione ad anelli (``daisy'') che rimuove la parte immaginaria
  a un loop è insufficiente a due loop (Bordag--Skalozub, conclusioni).
  Non esiste quindi un verdetto di stabilità perturbativo a questo
  ordine: il vuoto di Savvidy \textbf{non è stabile} contro le
  fluttuazioni, in accordo con la tesi di ``instabilità essenziale'' di
  Preparata.
\end{itemize}

\hypertarget{decomposizione-per-colori-e-modi-trasversi}{%
\subsection{4. Decomposizione per colori e modi
trasversi}\label{decomposizione-per-colori-e-modi-trasversi}}

Per settore (coefficienti esatti del nucleo di calore, termine
\(b^2 s\)):

\begin{longtable}[]{@{}lll@{}}
\toprule\noalign{}
settore & coefficiente & nota \\
\midrule\noalign{}
\endhead
\bottomrule\noalign{}
\endlastfoot
trasverso \(s_3 = +1\) & \(11/6\) & contiene il tachione \\
trasverso \(s_3 = -1\) & \(11/6\) & \\
coppia longitudinale \(s_3 = 0\) & \(-1/3\) & \\
fantasma & \(+1/3\) & cancella la coppia longitudinale \\
\end{longtable}

L'intero logaritmo a un loop è trasverso-paramagnetico: i due settori di
spin trasverso portano ciascuno \(11/6\), sommando all'\(11/3\) totale;
settore longitudinale e fantasma si cancellano esattamente.

\textbf{\(SU(3)\)} (fondo lungo \(\lambda_8\)): le radici cariche hanno
cariche \((1, \tfrac12, \tfrac12)\) in unità di \(gH\) --- i tre settori
tachiometrici di Cea. Ogni radice contribuisce come un vettore carico
complesso di \(SU(2)\) con la sua carica; il coefficiente del logaritmo
totale è \(\tfrac32 C\) con \(C = 11/(48\pi^2)\), più gli spostamenti
\(\ln q\) per radice.

\hypertarget{il-due-loop-verifica-e-tre-errata-documentate}{%
\subsection{5. Il due loop: verifica e tre errata
documentate}\label{il-due-loop-verifica-e-tre-errata-documentate}}

Dal calcolo delle definizioni del paper (loro eq. (3)) a \(T=0\),
\(\xi = 1\): il contributo a due loop è
\(W_2 = \tfrac32 g^2 B_2(0,b)^2\) con il ``tadpole'' magnetico
\[B_2(0,b) = \frac{b}{4\pi}\sum_{n,\sigma=\pm2}\int\frac{d^2k}{(2\pi)^2}\frac{1}{k^2 + b(2n+1+\sigma) - i0} = -\frac{b\,(\ln 2 - i\pi)}{16\pi^2},\]
esatto via funzione zeta: il polo e i termini \(\ln b\), \(\ln\mu\)
decadono identicamente perché
\(G(1) = 1 + \zeta(0,\tfrac12) + \zeta(0,\tfrac32) = 0\) (la
cancellazione coinvolge il tachione), e la somma residua
\(\sum \ln(2n+1+\sigma) = \ln 2 - i\pi\) segue dall'identità di Lerch
\(\zeta'(0,a) = \ln\Gamma(a) - \ln\sqrt{2\pi}\) (verificata
numericamente a \(10^{-30}\)).

\textbf{Errata candidati} (ciascuno fissato da un test):

\begin{itemize}
\tightlist
\item
  \textbf{F1} --- la loro eq. (58) stampa la correzione all'esponente
  \(3\ln^2 2/(98\pi^2)\) nella forma esponenziale e
  \(3\ln^2 2/(88\pi^2)\) nell'espansione. La minimizzazione della loro
  stessa (57) dà \(88\): il \(98\) è un refuso.
\item
  \textbf{F2} --- la loro (57) stampa la parte immaginaria a un loop
  \(-i b^2/(8\pi^2)\); due rotte indipendenti (e Nielsen--Olesen) danno
  \(b^2/(8\pi)\).
\item
  \textbf{F3} --- la loro (57) stampa il due loop
  \(g^2\ln^2 2\, b^2/(128\pi^4)\). Dalle loro stesse definizioni:
  \[\mathrm{Re}\,W_2 = \frac{3g^2 b^2 (\ln^2 2 - \pi^2)}{512\pi^4},\]
  cioè un fattore \(4/3\) sul coefficiente di \(\ln^2 2\) e una
  struttura \(-\pi^2\) assente nella forma stampata (differenza esatta:
  \(-g^2b^2(\ln^2 2 + 3\pi^2)/(512\pi^4)\)). Riprodotto
  indipendentemente dall'estrattore bibliografico della campagna.
\end{itemize}

\textbf{Minimo a due loop} (con il nostro \(W_2\)):
\(b_{\min} = \mu^2 \exp\left(-\frac{24\pi^2}{11g^2} - \frac{9(\ln^2 2 - \pi^2)g^2}{352\pi^2}\right)\).

\hypertarget{il-potenziale-in-funzione-di-blambda2}{%
\subsection{\texorpdfstring{6. Il potenziale in funzione di
\(b/\Lambda^2\)}{6. Il potenziale in funzione di b/\textbackslash Lambda\^{}2}}\label{il-potenziale-in-funzione-di-blambda2}}

Con running a un loop e miglioramento RG (\(\mu^2 = b\)),
\(x \equiv b/\Lambda^2\):
\[\frac{W(x)}{\Lambda^4} = \frac{11}{48\pi^2}\,x^2\left(\ln x - \frac12\right) + \kappa\,\frac{x^2}{\ln x},\qquad \kappa = \frac{9(\ln^2 2 - \pi^2)}{704\pi^2}\ \text{(derivato)},\quad \frac{3\ln^2 2}{176\pi^2}\ \text{(stampato)}.\]
A un loop il minimo è a \(x = 1\) esattamente: \(b_{\min} = \Lambda^2\);
\(W_{\min}/\Lambda^4 = -11/(96\pi^2)\).

\hypertarget{la-configurazione-classica-migliorata-il-vuoto-spaghetti}{%
\subsection{7. La configurazione classica migliorata: il vuoto
``spaghetti''}\label{la-configurazione-classica-migliorata-il-vuoto-spaghetti}}

Il tachione condensa in un reticolo bidimensionale di tubi di flusso
(Ambjørn--Olesen; rassegna Bordag eq. (146)--(152)). La densità del
condensato
\[\rho(z) = e^{-2\pi(\mathrm{Im}\,z)^2/\mathrm{Im}\,\tau}\,|\vartheta_3(\pi z\,|\,\tau)|^2\]
è esattamente periodica sotto le traslazioni magnetiche (verificato). Il
parametro di Abrikosov
\(\beta = \langle\rho^2\rangle/\langle\rho\rangle^2\), calcolato per
quadratura diretta dalla definizione: reticolo quadrato
\(\beta = 1.180340\), triangolare \(\beta = 1.159595\) (ottimale; la
mutazione obliqua è peggiore). L'energia classica diventa
\(E_{\rm cl} = \tfrac{H^2}{2}(1 - 1/\beta) < H^2/2\): il reticolo
\textbf{abbassa} l'energia classica del campo costante.

Riquantizzazione (approssimazione della letteratura: sommare il
potenziale a un loop del fondo omogeneo):
\[V(H) = \left(1-\frac1\beta\right)\frac{H^2}{2} + \frac{11}{48\pi^2}(gH)^2\left(\ln\frac{gH}{\mu^2} - \frac12\right),\qquad \ln\frac{gH_{\min}}{\mu^2} = -\left(1-\frac1\beta\right)\frac{24\pi^2}{11g^2},\]
e il minimo è più profondo di Savvidy del fattore
\(\exp\left(\frac{48\pi^2}{11 g^2\beta}\right)\) --- esponenzialmente
più profondo a accoppiamento debole.

Nota di convenzione (documentata, non risolta): la eq. (152) della
rassegna quota \(\vartheta_3(0, q_0) \simeq 1.2713\) per
\(q_0 = e^{-\pi}\); nella convenzione standard
\(\vartheta_3(0, e^{-\pi}) = 1.086434\), mentre
\(1.2713 = \vartheta_3(0, e^{-2})\). La loro eliminazione dei parametri
non è ricostruibile dal testo aperto; i nostri numeri sono derivati, non
trascritti.

\hypertarget{risultati-dichiarati}{%
\subsection{8. Risultati dichiarati}\label{risultati-dichiarati}}

\begin{enumerate}
\def\labelenumi{\arabic{enumi}.}
\tightlist
\item
  Il potenziale di Savvidy a un loop è confermato in ogni coefficiente:
  \(C = 11/(48\pi^2)\), minimo a \(b = \Lambda^2\),
  \(V_{\min} = -11\Lambda^4/(96\pi^2)\).
\item
  Il vuoto di Savvidy è \textbf{instabile}: tachione per ogni \(b>0\),
  \(\mathrm{Im}\,V \neq 0\) a uno e due loop; la tesi di Preparata
  dell'instabilità essenziale è confermata nel senso perturbativo
  accessibile (il suo testo completo non è accessibile: la conferma
  copre la struttura pubblicata da
  Savvidy/Nielsen--Olesen/Bordag--Skalozub, non le sue equazioni
  specifiche).
\item
  Il due loop pubblicato (Bordag--Skalozub) contiene tre errori (F1, F2,
  F3) con i valori corretti esatti sopra; la correzione F3 cambia il
  coefficiente del termine a due loop e l'esponente del minimo.
\item
  Il logaritmo a un loop è interamente trasverso-paramagnetico
  (\(2 \times 11/6\)); longitudinale e fantasma si cancellano.
\item
  La configurazione classica più vicina al vuoto quantistico è il
  reticolo triangolare di tubi di flusso (\(\beta = 1.159595\)), con
  energia classica ridotta del fattore \((1 - 1/\beta) = 0.1376\) e
  minimo esponenzialmente più profondo; la riquantizzazione a un loop
  attorno ad esso è calcolata sopra.
\item
  Il potenziale in funzione di \(b/\Lambda^2\) è la Sezione 6, con il
  termine a due loop corretto.
\end{enumerate}

\hypertarget{riferimenti}{%
\subsection{Riferimenti}\label{riferimenti}}

{[}1{]} G. Preparata, Nuovo Cim. A \textbf{96} (1986) 366. DOI:
10.1007/BF02833896 (abstract; testo paywalled). {[}2{]} G.K. Savvidy,
Phys. Lett. B \textbf{71} (1977) 133; riesposizione aperta: EPJC
\textbf{80} (2020) 165, arXiv:1910.00654. {[}3{]} N.K. Nielsen, P.
Olesen, Nucl. Phys. B \textbf{144} (1978) 376. {[}4{]} N.K. Nielsen, P.
Olesen, Phys. Lett. B \textbf{79} (1978) 304 (linee di vortice
elettriche). {[}5{]} J. Ambjørn, P. Olesen, Nucl. Phys. B \textbf{170}
(1980) 60 e 265 (condensato ``spaghetti''). {[}6{]} M. Bordag, V.
Skalozub, EPJC \textbf{82} (2022) 390, arXiv:2112.01043 (aperto).
{[}7{]} M. Bordag, Symmetry \textbf{15} (2023) 1137 (aperto). {[}8{]} P.
Cea, arXiv:2311.14791 (aperto). {[}9{]} Repository:
vantasnerdan/substrate-framework --- tracker \#44, issue \#45--\#50, PR
\#51--\#54; moduli \texttt{chromomagnetic\_background.py},
\texttt{chromomagnetic\_two\_loop.py},
\texttt{chromomagnetic\_sectors.py}, \texttt{spaghetti\_vacuum.py} con
50 test.

\emph{Fine del rapporto.}

\end{document}
